\problem{1.14}{Continuity equation and probability current}
\paragraph{Setup.} For a normalized one-dimensional wavefunction, derive the time rate of probability in $(a,b)$ and compute the current for Problem 1.9.
\paragraph{Derivation: continuity equation.}
Let $\rho=\Psi^*\Psi$. For a real potential,
\begin{align}
\partial_t\Psi&=\frac{\ii\hbar}{2m}\partial_x^2\Psi
-\frac{\ii}{\hbar}V\Psi,\\
\partial_t\Psi^*&=-\frac{\ii\hbar}{2m}\partial_x^2\Psi^*
+\frac{\ii}{\hbar}V\Psi^*.
\end{align}
By the product rule, the potential terms cancel:
\begin{align}
\partial_t\rho
&=\Psi^*\partial_t\Psi+\Psi\,\partial_t\Psi^*\\
&=\frac{\ii\hbar}{2m}\left(
\Psi^*\partial_x^2\Psi-\Psi\partial_x^2\Psi^*\right)\\
&=\frac{\ii\hbar}{2m}\partial_x\left(
\Psi^*\partial_x\Psi-\Psi\partial_x\Psi^*\right)
=-\partial_xJ,
\end{align}
where
\begin{equation}
\boxed{J(x,t)=\frac{\hbar}{2m\ii}
\left(\Psi^*\partial_x\Psi-\Psi\partial_x\Psi^*\right)
=\frac{\hbar}{m}\operatorname{Im}(\Psi^*\partial_x\Psi).}
\end{equation}
Define $P_{ab}(t)=\int_a^b\rho(x,t)\dd x$. Then
\begin{equation}
\boxed{\frac{\dd P_{ab}}{\dd t}
=-\int_a^b\partial_xJ\dd x=J(a,t)-J(b,t).}
\end{equation}
The sign has an intuitive meaning: inflow through the left boundary increases probability; outflow through the right boundary decreases it.
\paragraph{Dimensions.}
In one spatial dimension $[\Psi]=L^{-1/2}$ and
$[\Psi^*\partial_x\Psi]=L^{-2}$. Since $[\hbar/m]=L^2/T$,
\begin{equation}
\boxed{[J]=T^{-1},\quad \text{probability per unit time}.}
\end{equation}
\paragraph{Application to Problem 1.9.}
There $\Psi(x,t)=f(x)\e^{-\ii at}$ with a real Gaussian $f(x)$, so
$\Psi^*\partial_x\Psi=f(x)f'(x)$ is real. Hence
\begin{equation}
\boxed{J(x,t)=0\quad\text{for every }x,t.}
\end{equation}
\note{A stationary probability distribution can still carry a current in other systems, but the Gaussian ground state considered here does not.}

\problem{1.15}{Conservation of inner products}
\paragraph{Setup.} Show that the overlap of any two normalizable solutions of the same Schrödinger equation with real $V(x,t)$ is independent of time.
\paragraph{Derivation.}
Define the overlap
\begin{equation}
S(t)=\int_{-\infty}^{\infty}\Psi_1^*(x,t)\Psi_2(x,t)\dd x
=\langle\Psi_1(t)|\Psi_2(t)\rangle.
\end{equation}
Using $\partial_t|\Psi_2\rangle=-\ii\hat H|\Psi_2\rangle/\hbar$ and
$\partial_t\langle\Psi_1|=\ii\langle\Psi_1|\hat H/\hbar$ for a self-adjoint Hamiltonian,
\begin{align}
\frac{\dd S}{\dd t}
&=\frac{\ii}{\hbar}\langle\Psi_1|\hat H|\Psi_2\rangle
-\frac{\ii}{\hbar}\langle\Psi_1|\hat H|\Psi_2\rangle\\
&=\boxed{0}.
\end{align}
To see the boundary assumption explicitly, the difference in the two Hamiltonian matrix elements is
\begin{align}
\langle H\Psi_1|\Psi_2\rangle-\langle\Psi_1|H\Psi_2\rangle
&=-\frac{\hbar^2}{2m}\int
\left[(\partial_x^2\Psi_1)^*\Psi_2-
\Psi_1^*\partial_x^2\Psi_2\right]\dd x\\
&=-\frac{\hbar^2}{2m}
\left[(\partial_x\Psi_1)^*\Psi_2-\Psi_1^*\partial_x\Psi_2
\right]_{-\infty}^{\infty}=0,
\end{align}
provided the boundary term vanishes.
Thus
\begin{equation}
\boxed{\langle\Psi_1(t)|\Psi_2(t)\rangle
=\langle\Psi_1(0)|\Psi_2(0)\rangle.}
\end{equation}
\note{Taking $\Psi_1=\Psi_2$ recovers norm conservation. Taking initially orthogonal states preserves their orthogonality under common unitary evolution.}

\problem{1.16}{Parabolic compact-support wavefunction}
\paragraph{Setup.} For $\Psi(x,0)=A(a^2-x^2)$ on $[-a,a]$ and zero outside, compute normalization, expectations $\E x,\E p,\E{x^2},\E{p^2}$, and uncertainty product.
\paragraph{Derivation: normalization.}
Choosing $A>0$,
\begin{align}
1&=A^2\int_{-a}^a(a^2-x^2)^2\dd x\\
&=A^2\int_{-a}^a(a^4-2a^2x^2+x^4)\dd x\\
&=A^2\left[2a^5-\frac{4a^5}3+\frac{2a^5}5\right]
=A^2\frac{16a^5}{15}.
\end{align}
Thus
\begin{equation}
\boxed{A=\frac{\sqrt{15}}{4a^{5/2}}.}
\end{equation}
\paragraph{Derivation: position moments.}
The density is even, giving $\E x=0$. For the second moment,
\begin{align}
\E{x^2}
&=A^2\int_{-a}^a x^2(a^2-x^2)^2\dd x\\
&=A^2\cdot2\left[\frac{a^7}3-\frac{2a^7}5+\frac{a^7}7\right]\\
&=\frac{15}{16a^5}\frac{16a^7}{105}
=\boxed{\frac{a^2}{7}}.
\end{align}
Therefore $\boxed{\sigma_x=a/\sqrt7}$.
\paragraph{Derivation: momentum moments.}
Inside the support, $\partial_x\Psi=-2Ax$; the wavefunction is real and vanishes at $x=\pm a$. Thus
\begin{equation}
\E p=-\ii\hbar\int_{-a}^a\Psi\,\partial_x\Psi\,\dd x
=-\frac{\ii\hbar}{2}\left[\Psi(x,0)^2\right]_{-a}^a
=\boxed{0}.
\end{equation}
Integration by parts, with zero endpoint contribution, yields
\begin{align}
\E{p^2}&=-\hbar^2\int_{-a}^a\Psi\,\partial_x^2\Psi\,\dd x
=\hbar^2\int_{-a}^a|\partial_x\Psi|^2\dd x\\
&=4A^2\hbar^2\int_{-a}^a x^2\dd x
=4\frac{15}{16a^5}\hbar^2\frac{2a^3}{3}
=\boxed{\frac{5\hbar^2}{2a^2}}.
\end{align}
As a result
\begin{equation}
\boxed{\sigma_p=\sqrt{\frac52}\frac{\hbar}{a}},\qquad
\boxed{\sigma_x\sigma_p=\hbar\sqrt{\frac5{14}}
\simeq0.597614\hbar>\frac{\hbar}{2}.}
\end{equation}
\note{Although the derivative jumps at the support endpoints, the wavefunction itself is continuous and vanishes there, and the weak derivative is square-integrable. The integration-by-parts quadratic form above is legitimate.}

\problem{1.17}{A complex potential as a decay model}
\paragraph{Setup.} Let a phenomenological potential have a constant absorptive imaginary part,
$V(x)=V_0(x)-\ii\Gamma$ with $\Gamma>0$. Determine the total probability and lifetime.
\paragraph{Derivation.}
Taking the Schrödinger equation and its complex conjugate as in Problem 1.14, the real part $V_0$ cancels, but the imaginary contribution produces a sink:
\begin{equation}
\partial_t|\Psi|^2+\partial_xJ
=-\frac{2\Gamma}{\hbar}|\Psi|^2.
\end{equation}
Integrating over all space and assuming no boundary flux,
\begin{align}
P(t)&=\int_{\mathbb R}|\Psi(x,t)|^2\dd x,\\
\frac{\dd P}{\dd t}&=-[J]_{-\infty}^{\infty}
-\frac{2\Gamma}{\hbar}P
=\boxed{-\frac{2\Gamma}{\hbar}P}.
\end{align}
Separate variables and integrate:
\begin{equation}
\frac{\dd P}{P}=-\frac{2\Gamma}{\hbar}\dd t
\quad\Rightarrow\quad P(t)=P(0)\e^{-2\Gamma t/\hbar}.
\end{equation}
For $P(0)=1$, comparison with $P(t)=\e^{-t/\tau}$ yields
\begin{equation}
\boxed{\tau=\frac{\hbar}{2\Gamma}}.
\end{equation}
\note{A non-Hermitian effective Hamiltonian does not conserve the norm of the selected subsystem. The decreasing norm represents survival probability and does not contradict probability conservation in a larger closed system.}

\problem{1.18}{Thermal de Broglie wavelength and quantum regimes}
\paragraph{Setup.} Estimate the threshold for thermal de Broglie wavelength $\lambda$ to exceed a characteristic separation $d$ in a solid or ideal gas. Compare electrons, silicon nuclei, helium gas, and dilute atomic hydrogen in space.
\paragraph{Derivation: general condition.}
The stated thermal kinetic-energy estimate $p^2/(2m)=3k_BT/2$ gives a typical momentum $p=\sqrt{3mk_BT}$ and
\begin{equation}
\boxed{\lambda=\frac{h}{\sqrt{3mk_BT}}}.
\end{equation}
The criterion $\lambda\gtrsim d$ becomes
\begin{equation}
\boxed{T\lesssim T_{\rm q}(m,d):=\frac{h^2}{3mk_Bd^2}.}
\end{equation}
This estimate indicates when \emph{translational} wave effects are important; it is a rule of thumb, not a sharp phase transition.
\paragraph{Solids.}
For $d=0.3\,\mathrm{nm}=3.0\times10^{-10}\,\mathrm{m}$, using $h=6.62607015\times10^{-34}\,\mathrm{J\,s}$ and $k_B=1.380649\times10^{-23}\,\mathrm{J/K}$,
\begin{align}
T_{\rm q,e}&=
\frac{h^2}{3(9.10938\times10^{-31}\,\mathrm{kg})k_Bd^2}
\simeq\boxed{1.29\times10^5\,\mathrm K},\\
T_{\rm q,Si}&=
\frac{h^2}{3(28.085\,u)k_Bd^2}
\simeq\boxed{2.53\,\mathrm K},\qquad
u=1.66053906660\times10^{-27}\,\mathrm{kg}.
\end{align}
Thus conducting electrons are typically quantum-mechanical even at room temperature by this criterion, whereas lattice nuclei show comparable translational wave scales only at low temperature.
\paragraph{Ideal gases.}
For number density $n=N/V=P/(k_BT)$, estimate the spacing
\begin{equation}
d\simeq n^{-1/3}=\left(\frac{k_BT}{P}\right)^{1/3}.
\end{equation}
Substituting into $\lambda\gtrsim d$ and raising both sides to appropriate powers gives
\begin{align}
\frac{h^2}{3mk_BT}&\gtrsim\left(\frac{k_BT}{P}\right)^{2/3},\\
T^{5/3}&\lesssim\frac{h^2}{3mk_B}\left(\frac{P}{k_B}\right)^{2/3},
\end{align}
hence
\begin{equation}
\boxed{T\lesssim\frac1{k_B}\left(\frac{h^2}{3m}\right)^{3/5}P^{2/5}.}
\end{equation}
For helium atoms of mass $4.002602\,u$ at $P=101325\,\mathrm{Pa}$, this formula gives
\begin{equation}
\boxed{T_{\rm q,He}\simeq2.94\,\mathrm K.}
\end{equation}
At such temperatures ordinary helium may liquefy, so applying the \emph{ideal gas} estimate at atmospheric pressure is a formal scale comparison rather than a claim of a stable ideal-gas phase.
\paragraph{Hydrogen in interstellar conditions.}
For monatomic hydrogen $m\simeq1.00784u$, $T=3\,\mathrm K$,
\begin{equation}
\lambda=\frac{6.62607015\times10^{-34}}
{\sqrt{3(1.00784u)k_B(3\,\mathrm K)}}
\simeq\boxed{1.45\times10^{-9}\,\mathrm m}.
\end{equation}
Compared with $d\simeq1\,\mathrm{cm}=10^{-2}\,\mathrm m$,
\begin{equation}
\boxed{\frac{\lambda}{d}\simeq1.45\times10^{-7}\ll1}.
\end{equation}
Therefore this extremely dilute \emph{translational} gas is nondegenerate by the stated criterion; the internal electronic structure of its atoms is nevertheless quantum mechanical.
\note{The relevant distance scale differs across problems. A correct conclusion must specify whether the question concerns conduction electrons, atomic centers of mass, or internal bound states.}
